\section{Introduction}\label{sec:introduction}

The star operation permits a finite language expression to describe an
arbitrarily long repetition. Its nesting depth measures how many layers
of repetition an expression uses. Once complement is allowed, Boolean
conditions can replace some of those layers. The question studied here
is whether a fixed number of layers suffices for every regular language,
independently of its alphabet and recognizing automaton.

Fix a finite alphabet $\Sigma$. A \emph{generalized regular expression}
over $\Sigma$ uses $0$, $1$, the singleton letters of $\Sigma$, union,
concatenation, complement, and Kleene star. The constants denote
$\emptyset$ and $\{\eps\}$, and complement is always taken in $\Sigma^*$.
Its height is determined by
\[
\begin{aligned}
 \sht(0)=\sht(1)=\sht(a)&=0,\\
 \sht(P\cup Q)=\sht(PQ)&=\max\{\sht(P),\sht(Q)\},\\
 \sht(\neg P)&=\sht(P),&
 \sht(P^*)&=1+\sht(P).
\end{aligned}
\]
For a regular language $L\subseteq\Sigma^*$, write $h_\Sigma(L)$ for
the least height of an expression defining $L$. In particular,
complement preserves the nesting already present in its argument. Finite
Boolean combinations preserve a common height bound, and the full
language $\Sigma^*=\neg0$ has height zero.

\begin{theorem}\label{thm:main}
For every finite alphabet $\Sigma$ and every regular language
$L\subseteq\Sigma^*$,
\[
                            h_\Sigma(L)\le3.
\]
\end{theorem}

The expressions in the theorem use the letters of $\Sigma$ itself. Their
finite parameters and their lengths may depend on the language. The
construction gives an upper bound on nesting, without a bound on
expression size or running time.

\subsection{The problem and its mathematical context}

The corresponding \emph{ordinary} star-height problem omits complement.
Eggan connected ordinary star height with the cycle structure of
transition graphs, and Dejean and Sch\"utzenberger established unbounded
ordinary star height already over a two-letter
alphabet~\cite{Eggan1963,DejeanSchutzenberger1966}. These results concern
the operations allowed in ordinary expressions. Their lower bounds do
not carry over merely by adding complement.

Generalized height zero already has a different character: its languages
are the \emph{star-free} languages, which need not be finite. A finite
monoid $M$ recognizes a language when membership depends only on the image
of a morphism from $\Sigma^*$ to $M$. Sch\"utzenberger characterized the
star-free languages as exactly those recognized by finite aperiodic
monoids~\cite{Schutzenberger1965}. Here a monoid is aperiodic if every
element $m$ satisfies $m^r=m^{r+1}$ for some $r\ge1$. A simple language
requiring one star is $(aa)^*$ over $\{a\}$: its displayed expression has
height one, while every star-free unary language is finite or cofinite,
as induction under union, complement, and concatenation shows.

The algebraic approach has also produced substantial families of height
at most one. Pin, Straubing, and Th\'erien proved this bound for languages
recognized by finite nilpotent groups of class two, and for further
classes of finite monoids~\cite[Theorems~7.3 and~7.8]{PST1992}.
Another line of results concerns words in which a fixed nonempty factor
occurs contiguously a prescribed number of times modulo an integer
$n\ge2$, with overlapping occurrences counted. Bourne and Ru\v{s}kuc
proved height-one bounds for factors of length at most
three~\cite[Proposition~2.5]{BourneRuskuc2016};
Bourne subsequently extended the result to arbitrary fixed factors
in his thesis~\cite[Theorem~2.25]{Bourne2017}. These results explain how
complement can simplify nontrivial counting languages, but their
restricted families do not supply a bound for arbitrary regular languages.

A universal finite upper bound and a universal upper bound of one are
different assertions; Straubing distinguishes them explicitly in his
account of the problem~\cite[p.~537]{Straubing2002}.
Theorem~\ref{thm:main} gives a positive answer to the boundedness question.
It does not exhibit a language of generalized height greater than one,
and does not decide whether one always suffices.

Two companion constructions explain the setting of this proof. The paper
\emph{Finite Monoid Computations and a Uniform Generalized Star-Height
Bound} proves a bound of thirteen by finite computations, a prediction
argument, and a decision tree that retains the complete earlier
history~\cite[Theorem~1.1 and Sections~2--7]{HeightThirteen}.
\emph{Generalized Star Height at Most Four} proves a bound of four using
constancy on affine coordinate lines, several marker scales, and product
transport across moving periodic
boundaries~\cite[Theorem~1.1 and Sections~2--6]{HeightFour}.
The present construction obtains three through two split steps. Each of
the three papers gives its own complete proof; the numerical theorem of
one is not used to prove another. The prediction and multiscale methods
remain different constructions even though the bound here is smaller.

\subsection{Computing a word value with independent tests}

Every regular language has a recognizing morphism
$T:\Sigma^*\to M$ into a finite monoid. We therefore seek expressions
for the fibers of $T$. The construction parses a word into complete pieces called
\emph{episodes}, followed by a final remainder. The episode language $E$
is a prefix code, meaning that no member is a proper prefix of another.
This makes the episode boundaries unique.

An episode will act on a finite vector space by an affine map. To
propagate its state, we cut the episode and test the prefix for an input
state and the suffix for an output state. The tests are languages placed
in separate factors of a concatenation. Thus their auxiliary choices
are independent. A valid split has to meet two separate obligations:
every successful pair beginning at a true episode boundary must consume
that entire first episode and transmit its update; for every input state,
a successful pair must exist on each episode that is not assigned to a
specified skipped class.

Section~\ref{alg:section} proves a split identity that converts these
obligations into expressions for arbitrary episode sequences. If the
individual tests and the episode domain have height at most one and
sequence computations on the skipped class have height at most $H$, the
identity gives height at most $\max\{2,H+1\}$. The split and affine
recovery methods are developed in the two
companions~\cite[Lemmas~2.2--2.3]{HeightThirteen} and
\cite[Lemmas~2.1--2.2]{HeightFour}; all versions used here are proved
locally. An affine correction does not destroy the
original computation: the difference between the outputs on $x$ and on
zero recovers the linear part on $x$. Finite unions and intersections
express this subtraction without another star.

\subsection{Why two splits suffice}

The first split uses later cuts in an episode and a finite set of
\emph{tags}. A tag is a tuple of states, interval products, and indices
that a test proposes to use. The prefix and suffix choose their tags
separately. A clock restricts the pairs they can jointly realize: either
they agree, or there is only one possible unequal pair. In the second
case we define the episode's affine translation so that this one pair
transmits the prescribed input to the prescribed output. To use the
clock, however, both tests must calculate the same total time. A separate
end guard establishes that fact; it is not a consequence of the clock.

The clock also needs room to realize its permitted choices at actual
cuts. Outside an exceptional set of total times, every tag permits a
whole neighborhood of compatible prefix times, with a common positive
margin. The clock also supplies a finite list of candidate times for which
the number of exceptional pairs is bounded uniformly under translation
and sufficiently small simultaneous perturbations. These are the
additional clock properties used here;
Section~\ref{clock:section} compares the quantitative bounds.

When tags agree, the first split has two ways to work. If all the marker
boundaries are present, their successive intervals give a tuple of
monoid products. A finite shift table makes one coordinate of this tuple
irrelevant to the transmitted output. The prefix checks the earlier
products, the suffix checks the later products, and a finite certificate
checks every possible value of the product crossing their cut. If a
marker boundary is absent, a long periodic stretch instead allows the
cut to move through all clock residues while keeping its prefix product
fixed. This argument uses eventual periodicity of powers in a finite
monoid; it requires no cancellation.

The shift table is uniform in a stronger sense than a product-action
statement. One tuple length, chosen using only the finite entry alphabet and state
space, must work for every table from tuples to linear maps. A single correction table is chosen for each such
map, independently of the input vector. The coordinate that can be
omitted may depend on the tuple and vector. This quantifier order is
proved by a finite avoidance argument of the local-lemma type of
Erd\H{o}s and Lov\'asz~\cite{ErdosLovasz1975}.

The clock and marker conditions used to guarantee the first split's
completeness can fail on an intrinsic sublanguage $D\subseteq E$.
The second split handles every episode in $D$, with no further skipped
class. A suffix at an earlier cut can use the failed condition as a witness
to the actual origin. Those early cuts form
a fixed finite set, or \emph{stencil}, containing many copies of every
required tuple of metadata. The differences between distinct stencil
positions have different residues. Marker searches with widths indexed
by these residues make false witnesses sparse over all ordered pairs of
positions, including pairs with incorrect metadata.

End detection needs special care in this count. A smaller bounded search
fixes an episode's end when it finds an end marker; otherwise a periodic
seed is followed to its first mismatch. A larger bounded search supplies
the clock anchor whenever possible. If that larger search is absent at
one nearby proposed origin, all the smaller searches are absent, and
their overlapping seeds have one common mismatch. If every larger search
is present, all candidate clock anchors are already bounded data. In
either case the true origin's witness remains available before the suffix
requires uniqueness. After identifying it, the suffix must separately
check its exact endpoint. Counting over the whole stencil leaves a usable
copy of every metadata tuple and supplies the second split.

\subsection{From finite tests to height-three expressions}

The local marker construction is a finite-window instance of the
marker-versus-periodicity method~\cite[Lemma~2]{Krieger1982}; its full
priority construction and periodic-overlap argument are given in
Section~\ref{marker:section}. Section~\ref{clock:section} proves the
robust clock. Section~\ref{schedule:section} then chooses all finite
parameters in dependency order. In particular, a provisional rational
clock fixes the early stencil before the marker scale is chosen. A later
rational approximation preserves those fixed trials and makes the final
clock period coprime to all relevant short periods.

Sections~\ref{episode:section} and~\ref{outer:section} construct episodes,
force a common end under independently chosen roles, and prove the first
split. Section~\ref{decode:section} establishes the bound on all false
origins and the suffix decoder. Section~\ref{inner:section} uses the
remaining suffix to compute a total table indexed by every hypothetical
prefix product in $M$, including values that no prefix realizes. Applying
the shift-table lemma to this table gives the second split.

The last step, Section~\ref{sec:compilation}, controls expression height
for each component on its own domain. A standalone suffix, even with
false origin data, has its own exact-end test. It consequently lies in a
finite union of fixed-word templates $ac^ib$. The accepted exponents are
ultimately periodic, which gives height-one expressions over the original
letters. The incomplete remainder has the same kind of expression. The
inner split starts from empty skip computations and gives height two;
recovery makes those computations the skips for the outer split, giving
height three. Finite Boolean recovery and concatenation with the
remainder add no star. Tags, states, and residues serve only as finite
indices in these expressions.
